% Auto-generated: do not edit by hand.
% Sources: output/tables/{ardl,inflation}*_summary/sensitivity.csv,
% counterfactual_daily.csv, backtest_origins.csv, spec_robustness.csv,
% spec_range.csv, cf_bands_bootstrap.csv, event_study_coefs.csv,
% revision/{threshold_boundary_counts,price_endings_*,dispersion_summary}.csv

% ---------- D8/D9 tables ----------
\begin{table}[htbp]
\centering\small
\caption{Brand groups at the 28 October ceiling step (diesel)}
\label{tab:brand-28oct}
\begin{tabular}{lrcccc}
\toprule
 & & \multicolumn{2}{c}{27 Sep $\to$ 31 Oct} & \multicolumn{2}{c}{Below 2.212 ceiling} \\
\cmidrule(lr){3-4}\cmidrule(lr){5-6}
Group & Stations & $\Delta$ price (c/l) & Level 31 Oct & 28 Oct & 31 Oct \\
\midrule
Eni & 3\,741 & -0.0 & 2.190 & 99\% & 99\% \\
Q8 & 2\,584 & -3.8 & 2.216 & 2\% & 78\% \\
Api-Ip & 3\,385 & -2.7 & 2.229 & 5\% & 60\% \\
Esso & 1\,890 & -2.4 & 2.250 & 8\% & 32\% \\
Tamoil & 1\,374 & -3.4 & 2.240 & 3\% & 22\% \\
Other brands & 4\,032 & -1.6 & 2.271 & 9\% & 20\% \\
White pumps & 2\,584 & -1.8 & 2.286 & 10\% & 15\% \\
\bottomrule
\end{tabular}

\medskip
\footnotesize Mean gross diesel price change from 27 October (last day of
the 2.261 ceiling) to 31 October, and the share of observations strictly
below the mechanically adjusted 2.212 ceiling on 28 and 31 October. Eni's
network was already at the announced nominal level (2.19), hence compliant
under any ceiling in force.
\end{table}

\begin{table}[htbp]
\centering\small
\caption{Pesaran--Shin--Smith bounds tests for a level relationship
(net price, Brent), case 3, estimation window 2 Jan--27 Sep 2026}
\label{tab:ecm-bounds}
\begin{tabular}{lcccccc}
\toprule
 & Order & $F$ statistic & 5\% bounds $[I(0),I(1)]$ & $t$ statistic &
 ECT & Implied \\
 & & & & & (per day) & half-life \\
\midrule
Petrol & (2,8) & 1.40 & [4.92, 5.74] & -1.63 & -0.0134 & 52 days \\
Diesel & (1,8) & 0.92 & [4.92, 5.74] & -0.60 & -0.0047 & 148 days \\
\bottomrule
\end{tabular}

\medskip
\footnotesize $F$ and $t$ test the null of no level relationship between
the log net-of-tax price and log Brent in euro; critical bounds from
\textcite{pesaran2001}. ECT: coefficient on the lagged level in the
unrestricted error-correction representation. Both statistics lie far
below their $I(0)$ lower bounds: over this window the relationship is a
short-run one, and levels long-run elasticities are not identified ---
which is why the paper reports elasticity ranges and rests the
counterfactual on forecast validation.
\end{table}

\begin{table}[htbp]
\centering\small
\caption{Discrete-time logit hazard of first pricing below the ceiling}
\label{tab:hazard-dt}
\begin{tabular}{lccc}
\toprule
 & (1) Universe & (2) Excl.\ Eni + Q8 & (3) White pumps \\
\midrule
Pre-cap gap (10 c/l) & -0.888$^{***}$ & -1.143$^{***}$ & -1.551$^{***}$ \\
Nearest station (km) & -0.072$^{***}$ & -0.080$^{***}$ & -0.159$^{***}$ \\
Nearest day-0 adopter (km) & -0.040$^{***}$ & -0.053$^{***}$ & -0.051$^{***}$ \\
Population (1{,}000 per km\textsuperscript{2} cell) & -0.000 & +0.001 & +0.006 \\
Chain size (10 stations) & +0.004$^{***}$ & +0.003 & +0.057 \\
Diesel & +0.086$^{**}$ & +0.074$^{*}$ & -0.076 \\
\midrule
Day fixed effects & Yes & Yes & Yes \\
Brand fixed effects & Yes & Yes & -- (single brand) \\
Region fixed effects & Yes & Yes & Yes \\
\bottomrule
\end{tabular}

\medskip
\footnotesize One row per station--fuel unit per day at risk (463,594
unit-days; 32,974 adoptions); the event is the first day below the
ceiling. Log-odds; gestore-clustered standard errors suppressed for
readability (all starred coefficients significant at 0.1\% or better;
population insignificant everywhere). Day fixed effects absorb the
baseline hazard, which makes no proportional-hazards assumption
necessary; the estimated baseline spikes on the 1 October repricing
wave and the 6 October excise step.
\end{table}
